% A practitioner's handbook of correlated credibility.
% Statements and commentary only; every result links to its Lean proof.
% Built by CI (xelatex) and published next to the blueprint as handbook.pdf.
% Outline and open sections: handbook/outline.md.
\documentclass[11pt]{article}

\usepackage{amssymb, amsmath, amsthm}
\usepackage[margin=2.5cm]{geometry}
\usepackage{booktabs}
\usepackage[colorlinks=true, linkcolor=blue, urlcolor=blue]{hyperref}

\newtheorem{result}{Result}
\theoremstyle{remark}
\newtheorem*{caveat}{Watch out}

\newcommand{\hmu}{\hat\mu}
\newcommand{\Var}{\operatorname{Var}}
\newcommand{\cov}{\operatorname{cov}}
\newcommand{\bu}{\mathbf 1}
% \lean{Name}: link to the Lean proof (the site's find/ page redirects a
% declaration name to its line in the source).
\newcommand{\lean}[1]{\par\smallskip\noindent\textsf{\small Lean:
  \href{https://mikaelandblom-coder.github.io/credible-lean/find/\#doc/#1}{\detokenize{#1}}}\par}
\newcommand{\why}{\par\smallskip\noindent\textit{Why.}\ }
\newcommand{\use}{\par\smallskip\noindent\textit{Use it.}\ }

\title{Correlated credibility: a handbook\\
  {\large Useful results, each proved in Lean}}
\author{M.~Andblom}
\date{Draft}

\begin{document}
\maketitle

\begin{abstract}
  Credibility theory blends a segment's own estimate with the portfolio
  estimate. This handbook collects the results of one framework for doing
  that when segments are correlated, with classical Bühlmann--Straub as a
  special case. There are no proofs here: each result links to a
  machine-checked proof in Lean, and the
  \href{https://mikaelandblom-coder.github.io/credible-lean/}{blueprint}
  has the details. The comments say when to use a result, which way it
  pushes, and what can go wrong.
\end{abstract}

% ===========================================================================
\section{Setting}

\paragraph{The question.} A portfolio is split into cohorts
$j = 1, \dots, J$ (segments, years of underwriting, levels of a rating
factor), observed over periods $t = 1, \dots, T$. You want each cohort's
true key ratio $\mu_j$. You have two estimates: the cohort's own $\hmu_j$
(right target, noisy) and the portfolio's $\hmu$ (stable, but blind to
differences between cohorts). Credibility blends them:
\[
  \hmu_j^{\,b} = z_j\,\hmu_j + (1 - z_j)\,\hmu .
\]

\paragraph{The model.} Each observation is the overall level plus two
error terms:
\[
  \mu_{jt} = \underbrace{\mu + \tau\,\eta_j}_{\mu_j} + \sigma_{jt}\,\varepsilon_{jt},
  \qquad
  \cov(\varepsilon_{it}, \varepsilon_{js}) = \rho_{ij}\,\delta_{ts}.
\]
The \emph{cohort effect} $\eta_j$ says how far cohort $j$ really is from the
rest; $\tau^2$ is how much cohorts differ. The \emph{noise}
$\varepsilon_{jt}$ is what one period of data gets wrong; $\sigma_{jt}$ is
how much (small for a large cohort). Both are centred with unit variance,
and the cohort effects are uncorrelated with each other and with the noise.
Cohorts are linked only through their noise in the same period: a
storm or an inflation shock hits them all, with correlation $\rho_{ij}$.

Only first and second moments are assumed. No distribution is needed
anywhere in this handbook.

\paragraph{Why error terms?} Classical credibility gives each cohort a risk
parameter $\Theta_j$, drawn independently. Independence is built into that
language, so correlation between cohorts is awkward to state. Written as
error terms, it is one line: the formula for
$\cov(\varepsilon_{it}, \varepsilon_{js})$ above. This is the
random-effects notation of statistics.

\paragraph{Translation.} For readers used to Bühlmann--Straub
(Bühlmann \& Gisler, \emph{A Course in Credibility Theory}):

\begin{center}
\begin{tabular}{lll}
  \toprule
  Here & Bühlmann--Straub & Meaning \\
  \midrule
  $\mu$ & $m$ (collective premium) & overall level \\
  $\tau^2$ & $a$ & variance between cohorts \\
  $\sigma_{jt}^2$ & $\sigma^2/w_{jt}$ & noise variance (volume $w_{jt}$) \\
  $\rho_{ij}$ & $0$ & correlation of noise between cohorts \\
  $\sigma_{jt}$ & & a standard deviation, not a variance \\
  \bottomrule
\end{tabular}
\end{center}
The thesis this grew out of writes the cohort effect $\varepsilon_j$; here it
is $\eta_j$, to keep it apart from the noise.

\paragraph{The recipe.} Every result below follows the same three steps.
\begin{enumerate}
  \item Build $K$, the covariance matrix of the cohort estimates:
        $K_{ij} = \cov(\hmu_i, \hmu_j)$.
  \item Check that $K$ is positive definite.
  \item Read off the weights and the credibility factors (Section~\ref{sec:core}).
\end{enumerate}
Each model (correlated cohorts, unequal series lengths, a shared group level)
only changes step~1. Steps~2--3 are proved once.

% ===========================================================================
\section{The core: everything follows from \texorpdfstring{$K$}{K}}\label{sec:core}

\subsection*{The best blend}

\begin{result}[The best credibility factor]\label{res:bestZ}
  Let $g = \Var(\hmu_j - \hmu) > 0$. The error variance
  $\Var(z\hmu_j + (1-z)\hmu - \mu_j)$ is smallest at
  \[
    z^* = \frac{\Var(\hmu) - \cov(\hmu_j, \hmu) + \cov(\hmu_j, \mu_j) - \cov(\hmu, \mu_j)}{g},
  \]
  and only there.
  \lean{Credibility.BlendMoments.bestZ_optimal}
  \lean{Credibility.BlendMoments.bestZ_unique}
\end{result}
\use This holds for \emph{any} two unbiased estimates, not only the ones
below: plug in six second moments and you get the best $z$.
\why The error variance is a parabola in $z$ with curvature $g$; $z^*$ is
its vertex. If $g = 0$ the two estimates are the same and every $z$ is
equally good.

\subsection*{The best weights}

\begin{result}[Within a cohort: inverse-variance weights]\label{res:invVar}
  Among weights $c_t$ with $\sum_t c_t = 1$, the noise variance
  $\sum_t c_t^2\sigma_{jt}^2$ is smallest for
  \[
    a_{jt} = \frac{\sigma_{jt}^{-2}}{\sum_s \sigma_{js}^{-2}},
    \qquad\text{with minimum}\quad
    \Bigl(\sum_s \sigma_{js}^{-2}\Bigr)^{-1}.
  \]
  \lean{Credibility.invVarWeight_optimal}
  \lean{Credibility.invVarWeight_value}
\end{result}
\use $\hmu_j = \sum_t a_{jt}\mu_{jt}$. With $\sigma_{jt}^2 = \sigma^2/w_{jt}$
this is the volume-weighted mean, as in Bühlmann--Straub. Correlation
between cohorts does not change it: within a cohort, the noise of different
periods is uncorrelated.

\begin{result}[Across cohorts: GLS weights]\label{res:gls}
  Let $K$ be positive definite. Among weights $c_j$ with $\sum_j c_j = 1$,
  the portfolio variance $c^\top K c$ is smallest for
  \[
    b = \frac{K^{-1}\bu}{\bu^\top K^{-1}\bu},
    \qquad\text{with minimum}\quad
    s^2 = \frac{1}{\bu^\top K^{-1}\bu}.
  \]
  \lean{Credibility.glsWeight_optimal}
  \lean{Credibility.glsWeight_value}
\end{result}
\use $\hmu = \sum_j b_j\hmu_j$. With uncorrelated cohorts $K$ is diagonal
and $b_j \propto 1/K_{jj}$, the classical weights. With correlation, a
cohort that moves with many others gets \emph{less} weight: much of what it
says, the others already say. Weights can even turn negative.

\begin{result}[Every cohort covaries equally with the portfolio]\label{res:upsilon}
  If $K$ is positive definite, then $\cov(\hmu_j, \hmu) = s^2$ for every
  cohort $j$.
  \lean{Credibility.mulVec_glsWeight}
\end{result}
\why If some cohort covaried more with $\hmu$ than the others, shifting
weight away from it would lower $\Var(\hmu)$, and the weights would not be
optimal. This is what makes the next result simple.

\subsection*{The credibility factor}

\begin{result}[Credibility for correlated cohorts]\label{res:z}
  Let $K$ be positive definite, with $K_{jj} = \tau^2 + v_j$, where $v_j$ is
  the noise variance of $\hmu_j$. With $b$ and $s^2$ from
  Result~\ref{res:gls},
  \[
    z_j = \frac{(1 - b_j)\,\tau^2}{\tau^2 + v_j - s^2}
        = \frac{\cov(\hmu_j - \hmu,\; \mu_j)}{\Var(\hmu_j - \hmu)} .
  \]
  \lean{Credibility.Core.bestZ_eq}
\end{result}
\use This is the formula to implement. Compute $K$, invert it once, and
every $z_j$ follows. In the thesis this is eq.~(26).
\why The second form is a regression coefficient: $z_j$ answers ``how much
of the gap between my cohort and the portfolio is real?''. The numerator is
the real part of the gap: $\tau^2$, minus the share $b_j\tau^2$ of cohort
$j$'s own effect that already sits in $\hmu$. The denominator is the whole
gap, real plus noise.

\begin{result}[Uncorrelated cohorts: Bühlmann--Straub]\label{res:diag}
  If $K = \operatorname{diag}(\tau^2 + v_j)$ with positive entries and
  $J \ge 2$, then $b_j \propto 1/(\tau^2 + v_j)$ and
  \[
    z_j = \frac{\tau^2}{\tau^2 + v_j}.
  \]
  \lean{Credibility.Core.bestZ_of_diag}
\end{result}
\use In classical notation, $v_j = \sigma^2/w_j$ and this is
$z_j = w_j/(w_j + \sigma^2/a)$. So the correlated framework contains the
classical one exactly.

\subsection*{When the recipe applies}

\begin{result}[\texorpdfstring{$K$}{K} is positive definite]\label{res:posdef}
  In the model of Section~1, with $T \ge 1$ and the correlation table
  $(\rho_{ij})$ positive semidefinite, $K$ is positive definite. Its entries
  are $K_{ij} = \sigma_{ij}^2\rho_{ij} + \tau^2\delta_{ij}$, with
  $\sigma_{ij}^2 = \sum_t a_{it}a_{jt}\sigma_{it}\sigma_{jt}$.
  \lean{Credibility.Model.K_posDef}
  \lean{Credibility.Model.K_apply}
\end{result}
\begin{caveat}
  A covariance matrix is only positive \emph{semi}definite in general; here
  $\tau^2 > 0$ is what lifts it. If you estimate $\rho$ entry by entry, the
  table you get need not be positive semidefinite. Check it before inverting
  $K$; if it fails, repair it first (e.g.\ clip negative eigenvalues).
\end{caveat}

\begin{result}[A damped plug-in $\hat K$ is positive definite]\label{res:damped}
  Plug estimates $\hat\sigma_{jt}$, $\hat\rho_{ij}$ into $K$, multiply the
  off-diagonal correlations by a factor $\theta$ with $0 \le \theta < 1$, and
  use any $\hat\tau^2 \ge 0$. If the table $(\hat\rho_{ij})$ is positive
  semidefinite, then
  \[
    \hat K_{ij} = \hat\sigma^2_{ij}\,\theta^{[i\ne j]}\hat\rho_{ij} + \hat\tau^2\delta_{ij}
  \]
  is positive definite. This also holds with unequal series lengths.
  \lean{Credibility.Unequal.dampedK_posDef}
\end{result}
\use The estimate $\hat\tau^2$ is often negative; truncate it at zero and
damp the correlations a little (say $\theta = 0.95$). Then $b$, $s^2$ and every
$z_j$ are defined, whatever the data. Without damping, $\hat\tau^2 = 0$ and two
cohorts with $\hat\rho_{12} = 1$ can make $\hat K$ singular.
\why Damping moves $\hat K$ towards its diagonal, the Bühlmann--Straub $K$ of
uncorrelated cohorts, which is positive definite on its own. Damping does not
rescue an indefinite $\hat\rho$ for every $\theta$: that hypothesis stays.

\begin{caveat}
  You need at least two cohorts. With $J = 1$, $b_1 = 1$ and
  $s^2 = \tau^2 + v_1$, so Result~\ref{res:z} reads $0/0$: there is nothing
  to blend with.
\end{caveat}

% ===========================================================================
\section{Special cases}
\emph{To be written: two cohorts, credibility above one, equal covariances.}

\section{Estimating the parameters}
\emph{To be written.}

\section{The moments are real}
\emph{To be written.}

\section{Beyond the thesis}
\emph{To be written: unequal series lengths, merging cohorts, a shared group
level.}

\end{document}
